\documentclass[11pt,a4paper]{article}
\usepackage[margin=2.5cm,headheight=15pt]{geometry}
\usepackage{amsmath,amssymb,amsthm}
\usepackage{mathtools}
\usepackage{microtype}
\usepackage{hyperref}
\usepackage{cleveref}
\usepackage{enumitem}
\usepackage{mdframed}
\usepackage{xcolor}
\usepackage{titlesec}
\usepackage{booktabs}
\usepackage{fancyhdr}

\definecolor{navyblue}{RGB}{26,58,92}
\definecolor{midblue}{RGB}{44,95,138}
\definecolor{lightblue}{RGB}{232,240,248}
\definecolor{proofgray}{RGB}{242,242,242}
\definecolor{propcolor}{RGB}{240,248,240}
\definecolor{propborder}{RGB}{0,130,0}
\definecolor{opencolor}{RGB}{235,245,235}
\definecolor{openborder}{RGB}{0,120,0}
\definecolor{warncolor}{RGB}{255,245,220}
\definecolor{warnborder}{RGB}{180,100,0}

\hypersetup{colorlinks=true,linkcolor=midblue,citecolor=midblue,urlcolor=midblue}

\pagestyle{fancy}\fancyhf{}
\fancyhead[L]{\small\color{navyblue}\textit{Grouping Selection for the JGP+GSP Heuristic}}
\fancyhead[R]{\small\color{navyblue}\thepage}
\renewcommand{\headrulewidth}{0.4pt}

\titleformat{\section}{\large\bfseries\color{navyblue}}{}{0em}{}[\color{navyblue}\titlerule]
\titleformat{\subsection}{\normalsize\bfseries\color{midblue}}{}{0em}{}
\titleformat{\subsubsection}{\normalsize\itshape\color{midblue}}{}{0em}{}

\newmdenv[backgroundcolor=lightblue,linecolor=navyblue,linewidth=1.2pt,
  leftline=true,rightline=false,topline=false,bottomline=false,
  innerleftmargin=12pt,innerrightmargin=8pt,innertopmargin=8pt,
  innerbottommargin=8pt,skipabove=10pt,skipbelow=6pt]{thmbox}

\newmdenv[backgroundcolor=proofgray,linecolor=midblue,linewidth=0.8pt,
  leftline=true,rightline=false,topline=false,bottomline=false,
  innerleftmargin=10pt,innerrightmargin=8pt,innertopmargin=6pt,
  innerbottommargin=6pt,skipabove=6pt,skipbelow=6pt]{proofbox}

\newmdenv[backgroundcolor=propcolor,linecolor=propborder,linewidth=1.2pt,
  leftline=true,rightline=false,topline=false,bottomline=false,
  innerleftmargin=12pt,innerrightmargin=8pt,innertopmargin=8pt,
  innerbottommargin=8pt,skipabove=10pt,skipbelow=6pt]{propbox}

\newmdenv[backgroundcolor=opencolor,linecolor=openborder,linewidth=1.2pt,
  leftline=true,rightline=false,topline=false,bottomline=false,
  innerleftmargin=12pt,innerrightmargin=8pt,innertopmargin=8pt,
  innerbottommargin=8pt,skipabove=10pt,skipbelow=6pt]{openproblem}

\newmdenv[backgroundcolor=warncolor,linecolor=warnborder,linewidth=1.2pt,
  leftline=true,rightline=false,topline=false,bottomline=false,
  innerleftmargin=12pt,innerrightmargin=8pt,innertopmargin=8pt,
  innerbottommargin=8pt,skipabove=10pt,skipbelow=6pt]{warnbox}

\theoremstyle{plain}
\newtheorem{theorem}{Theorem}[section]
\newtheorem{lemma}[theorem]{Lemma}
\newtheorem{corollary}[theorem]{Corollary}
\theoremstyle{definition}
\newtheorem{definition}[theorem]{Definition}
\newtheorem{remark}[theorem]{Remark}
\newtheorem{conjecture}[theorem]{Conjecture}

\DeclareMathOperator{\OPT}{OPT}
\newcommand{\Kstar}{K^*}
\newcommand{\calG}{\mathcal{G}}
\newcommand{\ZSSP}{Z_{\mathrm{SSP}}^*}
\newcommand{\ZJGP}{Z_{\mathrm{JGP}}^*}

\setlist[itemize]{noitemsep,topsep=4pt}
\setlist[enumerate]{noitemsep,topsep=4pt}

% ─────────────────────────────────────────────────────────────────────────────
\begin{document}
% ─────────────────────────────────────────────────────────────────────────────

\begin{center}
  {\color{navyblue}\rule{\textwidth}{2pt}}\\[0.4cm]
  {\LARGE\bfseries\color{navyblue}The Job Sequencing and Tool Switching Problem}\\[0.3cm]
  {\large\itshape\color{midblue}Part VI: Grouping Selection for the JGP+GSP Heuristic}\\[0.2cm]
  {\color{navyblue}\rule{\textwidth}{2pt}}
\end{center}

\begin{mdframed}[backgroundcolor=lightblue,linecolor=navyblue,linewidth=1pt,
  innerleftmargin=14pt,innerrightmargin=14pt,innertopmargin=8pt,innerbottommargin=8pt]
\small\textbf{Scope and Status.}
This document addresses the following problem: given a pool of candidate groupings
(optimal and sub-optimal solutions to the JGP), how should we select groupings to
feed into the GSP so that the resulting SSP cost is minimised?
We present:
(i)~a proved characterisation of when a fixed grouping achieves the SSP optimum
via GTSP;
(ii)~the Max-Weight Hamiltonian Path (MWHP) result — \textbf{proved only for
instances where groups at distance 2 have disjoint tool sets};
(iii)~a Lagrangian relaxation framework for exploring the JGP/SSP tradeoff;
(iv)~ML-based grouping scoring approaches (proposed, not proved).
All proved vs.\ proposed results are clearly labelled.
\end{mdframed}

\vspace{0.5cm}
\tableofcontents
\newpage

% ─────────────────────────────────────────────────────────────────────────────
\section{Setup: The Grouping Selection Problem}
\label{sec:setup}
% ─────────────────────────────────────────────────────────────────────────────

After solving the JGP (via PORTA, Branch-and-Price, or a Golang enumeration),
we have access to:
\begin{itemize}
\item $\mathcal{G}^* = \{$ all JGP-optimal groupings with $\Kstar$ groups $\}$;
\item $\mathcal{G}_{>} = \{$ sub-optimal groupings with $K > \Kstar$ groups $\}$.
\end{itemize}
For each grouping $\calG = \{G_1,\ldots,G_K\}$, let $\OPT(\calG)$ denote the
optimal JGP+GSP heuristic cost when restricted to configurations from the
clusters $\{\mathcal{C}(G_k)\}_{k=1}^K$.

The goal is: find $\calG^{\dag} = \arg\min_{\calG \in \mathcal{G}^* \cup \mathcal{G}_{>}} \OPT(\calG)$.

Since $\OPT(\calG)$ is itself a GTSP and expensive to compute, we need efficient
\emph{proxies} to rank groupings without solving the full GTSP each time.

% ─────────────────────────────────────────────────────────────────────────────
\section{GTSP for a Fixed Grouping (Proved)}
\label{sec:gtsp_fixed}
% ─────────────────────────────────────────────────────────────────────────────

\begin{thmbox}
\begin{theorem}[GTSP equivalence for fixed grouping]
\label{thm:gtsp_fixed}
For a fixed grouping $\calG = \{G_1,\ldots,G_K\}$,
\[
  \OPT(\calG) = \text{min-cost Hamiltonian path on the GTSP over clusters }
  \{\mathcal{C}(G_k)\}_{k=1}^K
  \text{ with costs } d(C_i, C_j).
\]
\end{theorem}
\end{thmbox}

\begin{proofbox}
\begin{proof}
Any feasible SSP solution restricted to the grouping $\calG$ corresponds to:
selecting one configuration $C_k \in \mathcal{C}(G_k)$ per group, and choosing
an ordering of the groups. This is precisely a Hamiltonian path visiting one node
per cluster with arc costs $d(C_i,C_j)$, i.e., a GTSP (path variant).
Minimising over both configuration choice and ordering is the GTSP. \qed
\end{proof}
\end{proofbox}

\noindent\textbf{Practical solving.}
The Noon--Bean transformation \cite{Noon1993} converts a GTSP to a standard TSP
on $O(N)$ nodes ($N = \sum_k |\mathcal{C}(G_k)|$), solvable by Concorde or LKH-3.
For small $\Kstar$, Held--Karp DP for GTSP runs in $O(N^2 \cdot 2^{K^*})$.

Note that $|\mathcal{C}(G_k)| = \binom{m-|T^{(k)}|}{s(G_k)}$ is exponential in
the slack $s(G_k)$, so exact GTSP is tractable only for small instances or
small slack. Column generation can be used to implicitly represent the cluster.

% ─────────────────────────────────────────────────────────────────────────────
\section{Max-Weight Hamiltonian Path and Tool Overlap (Partly Proved)}
\label{sec:mwhp}
% ─────────────────────────────────────────────────────────────────────────────

\subsection{The Tool-Overlap Graph}

\begin{definition}[Tool-overlap graph]
For a grouping $\calG$, define the complete weighted graph $H_\mathrm{ov}$ on $K$
nodes (one per group) with edge weight:
\[
  w(G_k, G_{k'}) = \Bigl|\bigcup_{j\in G_k}T_j \;\cap\; \bigcup_{j\in G_{k'}}T_j\Bigr|.
\]
\end{definition}

Under \emph{one-directional} lookahead configurations --- each $C_k$ uses its slack to
pre-load tools of the \emph{next} group only (the per-transition device of Part~V,
\texttt{05\_jgp\_ssp\_gap\_analysis.tex}) --- the boundary cost
$k \to k+1$ is $r_k - w(G_k, G_{k+1})$ where $r_k = |\bigcup_{j \in G_k}T_j|$.
Minimising this one-directional cost is equivalent to maximising
$\sum_k w(G_{\pi(k)}, G_{\pi(k+1)})$ over orderings $\pi$: a
\textbf{Max-Weight Hamiltonian Path (MWHP)} on $H_\mathrm{ov}$. This is the optimal
ordering \emph{within the one-directional policy}; with positive slack, bidirectional
padding (a configuration retaining tools of \emph{both} neighbours) can be strictly
cheaper, so the MWHP value is in general an \emph{upper bound} on the true GSP optimum
(exact when slack is zero).

\subsection{The MWHP Result — with Proof Conditions}

\begin{thmbox}
\begin{theorem}[Optimal ordering as MWHP]
\label{thm:mwhp}
For a grouping $\calG$ where all groups have equal mandatory tool set size
$r = |\bigcup_{j \in G_k}T_j|$ for all $k$, and where groups at distance 2
have disjoint mandatory tool sets (i.e., $(\bigcup_{j \in G_k}T_j) \cap
(\bigcup_{j \in G_{k+2}}T_j) = \emptyset$ for all $k$), the optimal ordering
$\pi^*$ under lookahead configurations is a MWHP on $H_\mathrm{ov}$:
\[
  \OPT(\calG)_{\text{lookahead}} = (K-1)r \;-\;
  \max_\pi \sum_{k=1}^{K-1} w\!\left(G_{\pi(k)}, G_{\pi(k+1)}\right).
\]
\end{theorem}
\end{thmbox}

\begin{proofbox}
\begin{proof}
Under lookahead, $C_{\pi(k)} = T^{(\pi(k))} \cup X_{\pi(k)}$ where
$X_{\pi(k)} \subseteq T^{(\pi(k+1))} \setminus T^{(\pi(k))}$ and
$|X_{\pi(k)}| = s$ (the slack).
%% TODO-VERIFY(Claude-Fable 2026-06-10): missing hypothesis made explicit. For
%% X to exist we need |T^(k+1) \ T^(k)| = r - w >= s = b - r, i.e. consecutive
%% groups unmergeable (|T^(k) u T^(k+1)| > b). Automatic for JGP-optimal
%% groupings; NOT automatic for arbitrary groupings with K > K*.
Such an $X_{\pi(k)}$ exists because consecutive groups are unmergeable
($|T^{(\pi(k))}\cup T^{(\pi(k+1))}|>b$, automatic for JGP-optimal groupings and
assumed here), which gives
$|T^{(\pi(k+1))}\setminus T^{(\pi(k))}|=r-w\ge b-r=s$. For the last group
$\pi(K)$ there is no successor; fix $X_{\pi(K)}$ to be any tools \emph{outside}
$\bigcup_{k<K}T^{(\pi(k))}$ (the one-directional policy forbids backward
padding, which is what bidirectional padding would exploit). Then:
\[
  C_{\pi(k)} \setminus C_{\pi(k+1)}
  = T^{(\pi(k))} \setminus C_{\pi(k+1)},
\]
since $X_{\pi(k)} \subseteq T^{(\pi(k+1))} \subseteq C_{\pi(k+1)}$.
Now $C_{\pi(k+1)} = T^{(\pi(k+1))} \cup X_{\pi(k+1)}$ where
$X_{\pi(k+1)} \subseteq T^{(\pi(k+2))} \setminus T^{(\pi(k+1))}$.
The distance-2 disjointness condition ensures $T^{(\pi(k))} \cap X_{\pi(k+1)} = \emptyset$
(since $X_{\pi(k+1)} \subseteq T^{(\pi(k+2))}$ and $T^{(\pi(k))} \cap T^{(\pi(k+2))}=\emptyset$).
Therefore $T^{(\pi(k))} \cap C_{\pi(k+1)} = T^{(\pi(k))} \cap T^{(\pi(k+1))}$, and:
\[
  d\!\left(C_{\pi(k)}, C_{\pi(k+1)}\right)
  = \bigl|T^{(\pi(k))} \setminus C_{\pi(k+1)}\bigr|
  = r - \bigl|T^{(\pi(k))} \cap T^{(\pi(k+1))}\bigr|
  = r - w\!\left(G_{\pi(k)}, G_{\pi(k+1)}\right).
\]
Summing over $k=1,\ldots,K-1$ and minimising over $\pi$ is equivalent to
maximising $\sum_k w(G_{\pi(k)}, G_{\pi(k+1)})$. \qed
\end{proof}
\end{proofbox}

\begin{warnbox}
\textbf{Caveat on the distance-2 disjointness condition.}
The proof above uses $T^{(\pi(k))} \cap T^{(\pi(k+2))} = \emptyset$.
This holds for the sliding-window construction of
\cref{sec:sliding_window} (where groups are spaced far enough apart),
but does \emph{not} hold for general arbitrary groupings.
For general groupings, $T^{(\pi(k))} \cap X_{\pi(k+1)}$ may be non-zero,
and the formula $d(C_{\pi(k)}, C_{\pi(k+1)}) = r - w(G_{\pi(k)},G_{\pi(k+1)})$
no longer holds.
Using the MWHP heuristic for group ordering on general instances is a
reasonable approximation but is not guaranteed to find the true GSP optimum.
\end{warnbox}

\subsection{The Sliding-Window Instance Family}
\label{sec:sliding_window}

The MWHP ordering result applies to the following structured instance family, which
also realises the tool-counting lower bound of Part~V
(\texttt{05\_jgp\_ssp\_gap\_analysis.tex}).

\begin{definition}[Sliding-window instance $\mathcal{I}(K,r,\mathrm{ov},b)$]
%% TODO-VERIFY(Claude-Fable 2026-06-10): two fixes. (1) The JOBS were never
%% defined — added: one job per window. (2) Parameter range relaxed from r<b to
%% r<=b: the corollary below references the zero-slack case r=b, which the old
%% range excluded.
Parameters: $K$ groups, group tool set size $r$, overlap $\mathrm{ov}$
between adjacent groups, magazine capacity $b$, satisfying:
$0 \le \mathrm{ov} \le r/2$, $r \le b$, $2r - \mathrm{ov} > b$.
\textbf{Jobs:} one job per window, $T_{j_k}=T^{(k)}$ (so $|T_{j_k}|=r\le b$).
\begin{itemize}
\item \textbf{Tools:} $T = \{t_1,\ldots,t_m\}$ with $m = (K-1)(r-\mathrm{ov})+r$.
\item \textbf{Mandatory sets:} $T^{(k)} = \{t_{(k-1)(r-\mathrm{ov})+1},\ldots,
  t_{(k-1)(r-\mathrm{ov})+r}\}$ (consecutive windows of width $r$, shifted by
  $r-\mathrm{ov}$).
\item \textbf{Adjacency:} $|T^{(k)} \cap T^{(k+1)}| = \mathrm{ov}$;
  $|T^{(k)} \cup T^{(k+1)}| = 2r-\mathrm{ov} > b$ (adjacent groups cannot share
  a configuration $\Rightarrow$ they must be in different groups).
\item \textbf{Distance-2 disjointness:} By $\mathrm{ov} \le r/2$, the overlap
  region of $T^{(k)}$ and $T^{(k+1)}$ does not reach $T^{(k+2)}$, so
  $T^{(k)} \cap T^{(k+2)} = \emptyset$.
\end{itemize}
\end{definition}

For this family, $\Kstar = K$ (no adjacent groups can be merged) and the MWHP
result applies: the natural linear ordering $G_1 \to G_2 \to \cdots \to G_K$
maximises the overlap $\mathrm{ov}$ at each boundary and is the MWHP optimum, with
one-directional lookahead cost $(K-1)(r-\mathrm{ov})$.

\begin{corollary}[Optimal SSP cost for the sliding-window family]
\label{cor:sliding}
The one-directional MWHP schedule costs $(K-1)(r-\mathrm{ov})$, an \emph{upper} bound.
The true optimum is governed by the tool-counting bound of Part~V
(\texttt{05\_jgp\_ssp\_gap\_analysis.tex}): with $|U| = (K-1)(r-\mathrm{ov})+r$ tools,
\[
  Z_\mathrm{SSP}^{*} \;=\; |U|-b \;=\; (K-1)(r-\mathrm{ov}) - (b-r).
\]
Thus the one-directional MWHP value exceeds $Z_\mathrm{SSP}^{*}$ by exactly the slack
$b-r$, and the two coincide iff the groups are \emph{zero-slack} ($r=b$).
\end{corollary}
\begin{proofbox}
\begin{proof}
%% TODO-VERIFY(Claude-Fable 2026-06-10): proof was missing ("being met by a
%% bidirectional-padding schedule" was asserted without argument); added via the
%% Part V cost identity. Machine-checked on (K,r,ov,b) = (3,3,1,4), (4,3,1,4),
%% (3,4,2,5), (4,4,1,5): K*=K and Z*=|U|-b in all cases (/tmp run, ssp_verify).
$Z^{*}\ge|U|-b$ is the tool-counting bound. For achievability use the cost
identity of Part~V (cost $=(|U|-b)+R$ with $R$ the re-insertions): in the natural
window order every tool is required on a \emph{contiguous} set of jobs (each tool
lies in at most two adjacent windows by $\mathrm{ov}\le r/2$), so the lazy
schedule --- insert each tool at its first required job, evict only tools whose
last required job has passed --- never re-inserts: at the transition into job
$k{+}1$ exactly $r-\mathrm{ov}$ tools become dead and $r-\mathrm{ov}$ new tools
are needed, so capacity is never exceeded. Hence $R=0$ and
$Z^{*}=|U|-b$.
\end{proof}
\end{proofbox}

\begin{remark}[The MWHP excess is a policy artifact, not a heuristic gap]
%% TODO-VERIFY(Claude-Fable 2026-06-10): new observation from verification runs.
On all tested family members the \emph{full} GTSP optimum over the
$\Kstar$-grouping (bidirectional padding allowed) equals $Z^{*}$ --- e.g.\
$(K,r,\mathrm{ov},b)=(3,3,1,4)$: $Z^{*}=H=3$ while MWHP $=4$. So the
JGP$+$GSP heuristic has \emph{zero gap} on this family; only the restricted
one-directional policy pays the extra $b-r$. The family therefore does
\emph{not} witness heuristic gaps (consistent with the positive-gap witnesses of
Part~V having different job structure, e.g.\ the $8$-job sliding-window
\emph{ring}); it witnesses the value of bidirectional padding, and realises the
tool-counting bound exactly.
\end{remark}

\noindent\textbf{Practical use.}
For arbitrary instances, compute $H_\mathrm{ov}$ and solve MWHP as a proxy for
the optimal group ordering (tractable via Held--Karp DP in $O(K^2 2^K)$, feasible
for small $K^*$). Then run the full GTSP on the top-ranked orderings. This is
a heuristic shortcut, not an exact method, for general instances.

% ─────────────────────────────────────────────────────────────────────────────
\section{Lagrangian Relaxation: Exploring the JGP/SSP Tradeoff (Proposed)}
\label{sec:lagrangian}
% ─────────────────────────────────────────────────────────────────────────────

\begin{propbox}
\textbf{Status: Proposed direction, not yet proved or implemented.}
\end{propbox}

The JGP minimises group count $|\calG|$; the SSP minimises switching cost.
These objectives can conflict: a grouping with $K > \Kstar$ groups may yield
a lower switching cost than any $\Kstar$-group solution. This is \emph{proved} in
Part~V (\texttt{05\_jgp\_ssp\_gap\_analysis.tex}, Theorem on sub-optimal groupings):
on the $6$-job ring a $K{=}4$ grouping attains the SSP optimum $3$, strictly below the
JGP-optimal heuristic value $4$.

A Lagrangian penalty on extra groups:
\[
  L(\lambda) = \min_{\calG,\pi,\{C_k\}} \Bigl[
  \mathrm{cost}(\calG,\pi,\{C_k\}) + \lambda(|\calG| - \Kstar)
  \Bigr]
  \;=\;\min_{K\ge\Kstar}\bigl[z_K+\lambda(K-\Kstar)\bigr], \quad \lambda \ge 0,
\]
where $z_K$ is the best cost among $K$-group solutions (so $z_{\Kstar}=H$ and
$\min_K z_K=\ZSSP$ by the grouping-exactness theorem of Part~V).
%% TODO-VERIFY(Claude-Fable 2026-06-10): the earlier claims here were WRONG and
%% are replaced. (a) "By weak duality L(lambda) <= Z*_SSP for all lambda" is
%% FALSE: on the 6-ring (z_3=4, z_4=z_5=z_6=3), L(0.5)=3.5 and L(1)=4 > Z*=3.
%% (b) "Z*_SSP - L(lambda*) equals the LP integrality gap" was not meaningful:
%% in fact L(0) = Z*_SSP EXACTLY (grouping exactness), so max_lambda L >= Z*.
%% Correct facts below; numerically checked (/tmp/v06.py run, 2026-06-10).
The correct properties are:
\begin{itemize}
\item $L(0)=\ZSSP$ \emph{exactly} (grouping exactness, Part~V) --- the penalty-free
  problem \emph{is} the SSP, not a bound on it;
\item $L(\lambda)\le z_{\Kstar}=H$ for every $\lambda\ge0$ (plug in the
  $\Kstar$-group solution; the penalty vanishes), with $L(\lambda)\to H$ as
  $\lambda\to\infty$;
\item $L$ is concave, piecewise-linear, and nondecreasing in $\lambda$; its
  breakpoints are the slopes $(z_{\Kstar}-z_K)/(K-\Kstar)$, $K>\Kstar$, of the
  lower convex
  envelope of the Pareto set $\{(K,z_K)\}$, and the largest breakpoint is exactly
  the collapse threshold $\rho_c(I)$ of Part~VII
  (\texttt{07\_collapse\_variants.tex}, Theorem ``Collapse threshold formula'')
  with $\lambda=\rho$;
\item minimisers of $L(\lambda)$ for intermediate $\lambda$ are Pareto-efficient
  groupings with $K>\Kstar$ --- the useful output: a principled generator of
  sub-optimal groupings that may beat the $\Kstar$-restricted heuristic.
\end{itemize}
$L$ is thus a scalarisation knob for \emph{exploring} the (group count, switch
cost) trade-off --- the same parameterisation as the setup-time ratio $\rho$ of
Part~VII --- not a lower-bounding device for $\ZSSP$.

% ─────────────────────────────────────────────────────────────────────────────
\section{Chromatic Structure and the Conflict Graph (Proved Observation)}
\label{sec:chromatic}
% ─────────────────────────────────────────────────────────────────────────────

\begin{definition}[Job conflict graph \cite{Crama1994}]
$G_c = (J, E_c)$ where $(i,j) \in E_c \iff |T_i \cup T_j| > b$
(jobs $i$ and $j$ cannot share a configuration).
\end{definition}

Every feasible JGP grouping is a proper colouring of $G_c$ (no group contains a
conflicting pair), hence $\chi(G_c) \le \Kstar$. The converse \emph{fails}: a colour
class can be capacity-infeasible even with no internal conflict --- e.g.\ the three jobs
$\{t_1,t_2\},\{t_1,t_3\},\{t_1,t_4\}$ ($b=3$) are pairwise compatible but their union has
size $4$ --- so $\chi(G_c)=\Kstar$ holds only under a Helly-type condition on the tool
sets (it does hold for the ring instances). See Part~V
(\texttt{05\_jgp\_ssp\_gap\_analysis.tex}, \S\,Conflict-Graph Structure).

\begin{remark}[MWHP connects chromatic structure to switching cost]
Among all JGP-optimal groupings (each a proper $\Kstar$-colouring of $G_c$), the MWHP
criterion (\cref{sec:mwhp})
selects the colouring whose colour classes have maximum pairwise mandatory tool
overlap in the optimal ordering. This connects the chromatic structure of $G_c$
to the SSP objective. To our knowledge, this connection is not made explicit
in prior literature.
\end{remark}

% ─────────────────────────────────────────────────────────────────────────────
\section{Machine Learning Directions (Proposed)}
\label{sec:ml}
% ─────────────────────────────────────────────────────────────────────────────

\begin{propbox}
\textbf{Status: Proposed directions. None are proved or implemented yet.
Mathematical justifications are given where available.}
\end{propbox}

\subsection{GNN-Based Grouping Scorer}

\noindent\textbf{Goal.} Score a large pool of candidate groupings (from PORTA/Golang
enumeration) without solving the GTSP for each.

\noindent\textbf{Architecture.}
Encode the SSP instance as a bipartite job--tool graph $B = (J \cup T, E)$
with $(j,t) \in E \iff t \in T_j$.
A two-layer Graph Attention Network (GAT) produces job embeddings
$h_j \in \mathbb{R}^d$. For grouping $\calG$, represent group $G_k$ as the mean
embedding $\mathbf{g}_k = |G_k|^{-1}\sum_{j \in G_k} h_j$.
A second GNN on the $K$-node overlap graph $H_\mathrm{ov}$ (edge features
$w(G_k,G_{k'})$) predicts $\OPT(\calG)$.

\noindent\textbf{Mathematical justification.}
$\OPT(\calG)$ depends on the pairwise tool overlap between groups, which is
captured by $H_\mathrm{ov}$.
%% TODO-VERIFY(Claude-Fable 2026-06-10): corrected an overclaim. Xu et al. 2019
%% prove the OPPOSITE of "GNNs approximate any function on graphs": message-
%% passing GNN expressive power is BOUNDED by the 1-WL test. The honest
%% justification is permutation invariance + the fact that the relevant features
%% (pairwise overlaps) are given explicitly as edge weights of H_ov.
GNNs are permutation-equivariant, matching the permutation-invariance of the
SSP cost in the group labels; and since the decisive quantities (the pairwise
overlaps $w(G_k,G_{k'})$) are supplied \emph{explicitly} as edge features of
$H_\mathrm{ov}$, the predictor need not infer them --- sidestepping the
expressiveness limits of message-passing GNNs, which are bounded by the
1-Weisfeiler--Leman test \cite{Xu2019}. Caveat: $\OPT(\calG)$ also depends on
higher-order overlap patterns (e.g.\ bridge tools spanning three groups), which
pairwise features do not determine; this is precisely what the learned component
must capture, and no expressiveness guarantee is claimed.

\noindent\textbf{Training.} Solve instances exactly up to moderate size;
generate both $K=\Kstar$ and $K>\Kstar$ groupings; train on exact $\OPT(\calG)$ values.

\noindent\textbf{Guarantee.}
%% TODO-VERIFY(Claude-Fable 2026-06-10): bound corrected from epsilon to
%% 2*epsilon (estimation error hits both the selected and the best grouping),
%% and the role of the pool made explicit.
If the scorer approximates $\OPT(\calG)$ within $\pm\varepsilon$ on every
grouping in the candidate pool $\mathcal{P}$ (w.p.\ $\ge1-\delta$ jointly), the
grouping selected by the scorer has true cost
$\le \min_{\calG\in\mathcal{P}}\OPT(\calG) + 2\varepsilon$
(standard argmin-transfer argument). By grouping exactness (Part~V),
$\min_{\calG\in\mathcal{P}}\OPT(\calG)=\ZSSP$ \emph{provided} $\mathcal{P}$
contains an SSP-optimal grouping; otherwise the pool-coverage error adds on
top. No guarantee without the $\varepsilon$-accuracy assumption, which is
empirical.

\subsection{RL for Joint Grouping and Sequencing}

\noindent\textbf{Goal.}
Learn a policy that directly constructs groupings to minimise switching cost,
without being constrained to $\Kstar$ groups.

\noindent\textbf{MDP.}
State: partial job-to-group assignment.
Action: assign the next unassigned job to an existing feasible group or open a
new one (masked to satisfy $|\bigcup_{j\in G}T_j| \le b$).
Terminal reward: $R_T = -\mathrm{cost}(\calG, \pi^*, \{C_k^*\})$ (negative GSP cost,
computed after the episode ends).

\noindent\textbf{Training.}
REINFORCE with greedy rollout baseline \cite{Kool2019}:
$\nabla_\theta \mathcal{L} = -\mathbb{E}[(R_T - b)\nabla_\theta \log \pi_\theta]$.

\noindent\textbf{Key distinction.}
A JGP solver minimises group count. This RL policy minimises switching cost
directly, and may discover that $K > \Kstar$ groups are beneficial --- the
sub-optimal-grouping phenomenon proved in Part~V
(\texttt{05\_jgp\_ssp\_gap\_analysis.tex}).

\noindent\textbf{Guarantee.}
No worst-case bound; distributional convergence only.

\subsection{Learning-Augmented Column Generation}

\noindent\textbf{Setup.}
The JGP is solved via column generation \cite{Crama1994Column}. The pricing
subproblem (finding a maximal feasible group with minimum reduced cost) is
itself NP-hard. An ML predictor $P_\phi : (u, B) \to G'$ (where $u$ are dual
variables and $B$ is the instance graph) acts as a fast heuristic oracle.

\noindent\textbf{Protocol.}
At each CG iteration: query $P_\phi$; check reduced cost of $G'$ (cheap); if
negative, add to master; else fall back to exact pricing.
All ML-suggested columns are verified before use; LP optimality is certified
by the fallback. Exact LP optimality is preserved; ML only reduces pricing
calls. Analogous approaches substantially reduce pricing effort on related
problems \cite{Morabit2021}.
%% TODO-VERIFY(Claude-Fable 2026-06-10): the earlier "50--80%" figure could not
%% be verified against Morabit et al. and is removed pending a check of the
%% paper's actual reported reductions; reinstate with a page/table reference.

% ─────────────────────────────────────────────────────────────────────────────
\section{Summary: What Is Proved vs.\ Proposed}
\label{sec:summary}
% ─────────────────────────────────────────────────────────────────────────────

\begin{center}
\begin{tabular}{lll}
\toprule
Result & Status & Location \\
\midrule
GTSP for fixed grouping & \textbf{Proved} & \cref{thm:gtsp_fixed} \\
MWHP = optimal ordering (distance-2 disjoint instances) & \textbf{Proved} & \cref{thm:mwhp} \\
MWHP for general instances & \textbf{Open / Approximation only} & \cref{sec:mwhp} \\
Sliding-window family: $\ZSSP = |U|-b$ & \textbf{Proved} & \cref{cor:sliding} \\
Lagrangian Pareto frontier & \textbf{Proposed} & \cref{sec:lagrangian} \\
Chromatic/MWHP connection & \textbf{Proved observation} & \cref{sec:chromatic} \\
GNN grouping scorer & \textbf{Proposed} & \cref{sec:ml} \\
RL for joint grouping & \textbf{Proposed} & \cref{sec:ml} \\
ML-augmented CG & \textbf{Proposed} & \cref{sec:ml} \\
\bottomrule
\end{tabular}
\end{center}

% ─────────────────────────────────────────────────────────────────────────────
\bibliographystyle{abbrvnat}
\bibliography{references}
% ─────────────────────────────────────────────────────────────────────────────
\end{document}
